% Ad Familiares 13.71 --- headnote
% ad-familiares-13-71, 46 BC, Rome

\textit{Cicero at Rome to P. Servilius Isauricus, proconsul of
Asia, written in the course of 46 BC (the manuscript dateline:
\textit{Scr. Romae, ut videtur, a. 708}). A recommendation of
T. Agusius, framed at greater warmth than the routine cards:
Agusius had been the companion of Cicero's exile in 58--57 BC,
``the partner of all my journeyings, voyages, labours, and
perils,'' and is recommended to Servilius ``as one of my
household and most intimate connections.''}

\textit{The closing piece in the Servilius recommendation cluster
preserved here at the end of \textit{Ad Familiares} 13. The
opening admission --- that Cicero is bound to recommend many,
since the closeness of his connection with Servilius is public
knowledge, but that the cases differ --- is the standard
discriminator of the genre, and exists to tell the recipient that
this particular recommendation is to be taken at the higher
weight. The remembrance of the exile period --- which still
defines for Cicero the bond he owes those who shared it --- is the
characteristic touch.}
